% Transcription — fonds Alexandre Grothendieck, shelfmark 29, batch 9 (pages 161-180).
% Established from the facsimile archives/batches/29/batch-09.pdf (a mirror of
% https://grothendieck.umontpellier.fr/29.pdf), per the skill
% transcribe-grothendieck.
%
% Pass: Opus 5 (claude-opus-5), 2026-08-29 — first pass, unchecked against the
% pages by a human.
%
% Sixteen of the twenty pages are transcribed. Four are not:
%
%   - page 163, an otherwise blank sheet inscribed top right, in his hand,
%     « <struck>Le Groupe</struck> / Progroupe fondamental / pour les
%     revêtements infinis principaux »;
%   - page 172, the same, inscribed « Plan »;
%   - page 174, the same, inscribed « Π₂ »;
%   - page 177, the same, inscribed « π₀, π₁ et formalisme simplicial
%     (Revêtements infinis) ».
%
% All four are wrappers naming the run that follows them, exactly as page 143
% of batch 8, page 123 of batch 7 and page 111 of batch 6; they get no \page{},
% and their words are carried into the section headings instead, where they are
% marked as his.
%
% The batch falls into six runs, and they are legible in very unequal degrees.
%
%   Pages 161-162 — the image of π₁(X) in π₁(Y) for f : X → Y proper dominant
%     with Y normal, and the nodal counter-example that shows normality cannot
%     be dropped. Fluent hand, mostly readable.
%   Pages 164-169 — the run the wrapper of page 163 names: the functor
%     π ↦ π¹(S,ξ;π) on principal homogeneous fibrations, its left exactness,
%     the descending filtration of structural groups one may restrict to, and
%     the strict pro-representability that follows. This is the hardest hand in
%     the folder so far: the formulas carry, the connective prose very often
%     does not, and the apparatus says so rather than inventing. The lower
%     third of page 166 is struck through with diagonal strokes and nothing in
%     it can be read.
%   Pages 170-171 — H¹_C(S,G) = lim H¹(T/S,G) for a discrete group G, the five
%     classes C₁ … C₅ of covering families, and the comparisons over a Dedekind
%     ring, a complete discrete valuation ring, a complete local ring. Written
%     out fair and almost entirely legible.
%   Page 173 — the plan in seven sections for the two base-change theorems.
%   Pages 175-176 — π₂ : the Hochschild–Serre computation of π²(X,G) against
%     the tower of coverings, and the exact sequence its edge maps give.
%   Pages 178-180 — the proposition on strict pro-representability of
%     G ↦ π¹(S'/S,ξ;G) for a submersive S' → S, the note recording what he does
%     not know about effectivity, and the simplicial diagram translating
%     descent data on E_{X₀} into cocycles on π₀.
%
% The dating is the inventory's own, for the whole shelfmark.
\documentclass[11pt,a4paper]{article}
% Le préambule commun à toutes les transcriptions du fonds Grothendieck.
%
% Il tient en une idée : **l'incertitude se compose, elle ne se résout pas.**
% Une transcription qui lisse un mot douteux a détruit la seule chose pour
% laquelle on la faisait — un lecteur doit pouvoir savoir, à chaque endroit,
% ce qui était sur le feuillet et ce qui a été ajouté. Les six macros qui
% suivent sont donc l'appareil critique, et non de la décoration.
%
% Les mêmes commandes sont reconnues par `scripts/render.mjs`, qui produit la
% vue de lecture du site. Une macro ajoutée ici doit l'être aussi là-bas,
% faute de quoi le rendu échouera bruyamment — ce qui est voulu : un rendu
% silencieusement faux est pire qu'un rendu absent.

\ProvidesPackage{grothendieck}[2026/08/08 Transcriptions du fonds Grothendieck]

\usepackage{amsmath,amssymb,amsthm}
\usepackage{mathrsfs}
\usepackage{stmaryrd}
% Les diagrammes commutatifs sont partout dans ce fonds : c'est la langue
% même de la géométrie algébrique de Grothendieck, et une transcription qui
% les rendrait en prose ne transcrirait rien.
\usepackage{tikz-cd}
% Français par défaut — c'est la langue des feuillets — mais l'anglais est
% chargé aussi : la traduction et le résumé sont des documents anglais, et la
% typographie française (espaces avant les deux-points) y serait une faute.
% Ces documents basculent par \selectlanguage{english} en tête de corps.
\usepackage[english,french]{babel}
\usepackage{fontspec}
\usepackage{geometry}
\geometry{a4paper,margin=25mm,marginparwidth=28mm}
\usepackage{marginnote}
\usepackage{xcolor}
% ulem, not soul. `soul`'s \st and \ul are built on a character-by-character
% scan that XeLaTeX's font machinery breaks: the document fails with an
% undefined control sequence rather than degrading. ulem does the same two jobs
% and is Unicode-engine safe. `normalem` stops it hijacking \emph, which the
% transcriptions use for Grothendieck's own emphasis.
\usepackage[normalem]{ulem}
\usepackage{hyperref}
\hypersetup{colorlinks=true,linkcolor=black,urlcolor=black,pdfborder={0 0 0}}

% --- Métadonnées du lot -----------------------------------------------------
% Reprises telles quelles par le script de rendu, qui les affiche en tête de
% la vue de lecture. Elles fixent ce que le fichier prétend couvrir : sans
% elles, une transcription flotte sans point d'attache dans un fonds de seize
% mille feuillets.

\newcommand{\folder}[1]{\gdef\grothFolder{#1}}
\newcommand{\batch}[1]{\gdef\grothBatch{#1}}
\newcommand{\leaves}[2]{\gdef\grothFirst{#1}\gdef\grothLast{#2}}
\newcommand{\foldertitle}[1]{\gdef\grothFoldertitle{#1}}
\gdef\grothFolder{?}\gdef\grothBatch{?}\gdef\grothFirst{?}\gdef\grothLast{?}\gdef\grothFoldertitle{}

% --- Le résumé --------------------------------------------------------------
%
% Il ouvre la lecture modernisée, sous le titre « Résumé », et il est composé
% en petit. Ce n'est pas un ornement : c'est le seul passage du document qui
% ne soit pas des mathématiques, et un lecteur doit pouvoir le reconnaître
% avant de l'avoir lu — puis le sauter s'il connaît déjà le sujet, ou s'y
% arrêter s'il ne le connaît pas. Le filet qui le suit marque la reprise du
% corps.

\newenvironment{resume}
  {\small\setlength{\parskip}{0.45em}}
  {\par\medskip\hrule\medskip}

% --- Mots-clés ---------------------------------------------------------------
%
% En anglais, en clôture du résumé de la lecture modernisée : le vocabulaire
% moderne sous lequel chercher ce que les feuillets construisent. Le manifeste
% du site (`npm run manifest`) les extrait pour étiqueter la cote entière —
% cette ligne est la source unique des tags, il n'y a pas de fichier à côté.
\newcommand{\keywords}[1]{\par\smallskip\noindent
  {\footnotesize\textsc{Keywords} --- \textit{#1}}\par}

% --- L'appareil critique ----------------------------------------------------

% Le numéro de feuillet, en marge. C'est l'ancre qui permet de régler un
% désaccord sur une formule en regardant *un* feuillet plutôt qu'une chemise
% de six cents. Le site s'en sert aussi pour tourner les pages du fac-similé
% quand on fait défiler la transcription.
\newcommand{\leaf}[1]{%
  \marginnote{\footnotesize\color{gray!70}\textsf{#1}}%
  \label{leaf:#1}\ignorespaces}

% Illisible : on ne devine pas. Un mot inventé qui se lit comme les autres est
% le pire résultat possible.
\newcommand{\ill}{\textcolor{red!60!black}{[\,\ldots\,]}}

% Lecture douteuse : le mot est donné, et signalé comme incertain.
\newcommand{\uncertain}[1]{\uline{#1}}

% Ajout de l'éditeur — une lettre manquante, un mot sous-entendu.
\newcommand{\add}[1]{\textcolor{blue!50!black}{[#1]}}

% Biffé par Grothendieck lui-même. Ce qu'il a rayé fait partie du document :
% une rature dit souvent où la pensée a changé de direction.
\newcommand{\struck}[1]{\sout{#1}}

% Note du transcripteur, sur l'état matériel du feuillet.
\newcommand{\note}[1]{\par\smallskip{\small\color{gray!60!black}\textit{[#1]}}\par\smallskip}

% Note marginale de Grothendieck, distincte du corps du texte.
\newcommand{\marginal}[1]{\par\smallskip\noindent\hspace{1em}%
  {\small\color{gray!60!black}#1}\par\smallskip}

% --- Titre ------------------------------------------------------------------

\AtBeginDocument{%
  \begin{center}
    {\large\bfseries Fonds Alexandre Grothendieck --- cote n\textsuperscript{o}~\grothFolder}\\[2pt]
    {\itshape\grothFoldertitle}\\[4pt]
    {\small feuillets \grothFirst--\grothLast\ (lot \grothBatch)}\\[2pt]
    {\footnotesize\color{gray!60!black}Transcription établie d'après le fac-similé numérisé
     par l'Université de Montpellier.\\
     Les lectures incertaines sont soulignées, les illisibles notées [\,\ldots\,],
     les ajouts entre crochets.}
  \end{center}
  \vspace{1em}}

\endinput


\folder{29}
\batch{9}
\pages{161}{180}
\dating{1959-1969}
\watermark{Édition de démonstration}
\foldertitle{Groupe fondamental [Autour de SGA 4 et SGA 7] : lettres (1959, 1967, 1969), tapuscrits et copies de tapuscrit annotés (s.d.), notes manuscrites (s.d.)}

\begin{document}

\section*{L'image de $\pi_1(X)$ dans $\pi_1(Y)$, et le rôle de la normalité (pages 161 et 162)}

\page{161}

Soit $f : X \to Y$ un morphisme \marginal{propre} \uncertain{dominant}, $Y$
\emph{normal} connexe \marginal{(loc. noeth.)}.

Alors $\pi_1(X) \to \pi_1(Y)$ a un \uncertain{conoyau} fini (ou plutôt,
\ill{} le cas \uncertain{où} $f$ est \uncertain{génériquement} \ill{} sur $Y$,
et \uncertain{montre} que $\pi_1(\uncertain{y}) \to \pi_1(Y)$ est
\uncertain{surjectif}, --- on est ramené au cas où $Y = \mathrm{Spec}(K)$.

On voit \ill{} qu'on peut supposer $X$ \ill{} sur $k$, et
\ill{} [remplaçant $K$ par une extension \uncertain{séparable} finie
convenable $K'$, et \uncertain{désignant} \uncertain{par} $X'$ \ill{}
\uncertain{ses composantes irréductibles} \ill{}, \ill{} un
\uncertain{composant connexe} de $X'$]. Mais alors $\pi_1(X) \to \pi_1(Y)$ est
\ill{}

\emph{N.B.} Si $Y$ n'est pas \emph{normal}, le th. est faux,
\struck{si} p. ex. $X$ est le normalisé de $Y$. Ex. $Y$ courbe
\uncertain{sur un} $k$ \uncertain{alg. clos}, avec pt double ordinaire, $X'$
son normalisé. Alors on a $\pi_1(X) = 0$, $\pi_1(Y) = \widehat{\mathbb{Z}}$
\uncertain{privé de sa $p$-composante}.

On voit \uncertain{facilement} que dans le \uncertain{dernier} cas \ill{}, on
ne peut pas former $\underline{\mathrm{Hom}}_Y(X, Y \sqcup Y)$ [car si
\struck{un pt} $\overline{y}$ \ill{} de $Y$, et $y'$, $y''$ deux pts
géométriques de \ill{} au-dessus, alors qu'il \ill{}].

Supposons maintenant \ill{} $f : X \to Y$ \emph{universellement ouvert}
[p. ex. $Y$ normale] je crois qu'on

\page{162}

pour \struck{des} \ill{} contenus dans $X_1$ étale sur $X$.
P.-ê. plutôt :

\begin{tikzcd}[column sep=small, row sep=small]
& X \arrow[d] \\
S & T \arrow[l]
\end{tikzcd}

\note{le diagramme est encadré sur la page et porte au-dessus de lui la
mention $f_{*}(X_1)$ ; c'est tout ce que la page ajoute}

\section*{Progroupe fondamental pour les revêtements infinis principaux : le foncteur $\pi \mapsto \pi^1(S,\xi;\pi)$ (pages 164 à 169)}

\note{le feuillet qui ouvre la suite --- page 163, par ailleurs nu --- porte de
sa main, en haut à droite : « \struck{Le Groupe} Progroupe fondamental pour les
revêtements infinis principaux ». Le titre de cette section en vient}

\page{164}

Le $\widetilde{\pi}_1(S,\xi)$, $S$ préschéma connexe, $\xi$ pt géom.

Pour tout groupe $\pi$, soit
\[
  \pi^1(S,\xi;\pi)
\]
l'ens. \ill{}, à isom. près, de \struck{esp.} \marginal{fibrés} principaux
homogènes sur $S$, de groupe structural $\pi$, \uncertain{ponctués} au dessus
de $\xi$.

\begin{itemize}
\item[$\alpha$)] $\pi \rightsquigarrow \pi^1(S,\xi;\pi)$ est \emph{exact à
gauche}.
\end{itemize}

En effet

\begin{itemize}
\item[a)] Commute aux produits finis --- c'est trivial.
\item[b)] Transforme \uncertain{une injection} $\pi \to \pi'$ en
\uncertain{une injection}.
\end{itemize}

En effet, \ill{} : \ill{} \ill{} un isom.
$P_1' \xrightarrow{\ u'\ } P_2'$ de deux $\pi'$-fibrés \ill{}, il existe
\ill{} un et un seul tel isom. qui \ill{} applique
$\xi_1' \mapsto \xi_2'$. Si donc $P_1'$ et $P_2'$ sont \uncertain{munis} de
$P_1$, $P_2$ deux
\[
  \begin{cases} P_i \subset P_i' \\ \xi_i = \xi_i' \end{cases}
\]
\ill{} \struck{applique} \ill{} $u : P_1 \to P_2$. En effet,
\struck{\ill{}}

\page{165}

\struck{\ill{}} $P_i$ se déduit de $P_i'$ par une \emph{section} de
$P_i'/\pi$ \uncertain{pointée} par $\xi_i'$. Si elle \uncertain{existe}, il ne
peut y en avoir qu'une seule telle section.

\begin{itemize}
\item[c)] Si $\pi \to \pi' \rightrightarrows \pi''$ est \emph{exacte}, alors
\struck{\ill{}} $\varphi(\pi) \to \varphi(\pi') \rightrightarrows
\varphi(\pi'')$ \uncertain{aussi}.
\end{itemize}

Je \struck{veux} \uncertain{vais} \ill{} que si $P'$ est \ill{} $\pi'$,
\ill{} tel que \ill{} les deux \ill{} est \uncertain{muni} de $\xi'$,
\ill{} un isom.
\[
  P' \times_{\pi_1}(\pi'', \varphi_1) \xrightarrow{\ u\ }
  P' \times_{\pi_1}(\pi'', \varphi_2)
\]
alors on peut \ill{} le groupe structural. Or soit $Q$ le
\uncertain{sous-schéma} de $P$ des \uncertain{coïncidences} de $u_1$ et $u_2$.

\marginal{$\pi \to \pi' \rightrightarrows \pi''$, et $P' \xrightarrow{i_1}
P_1''$, $P' \xrightarrow{i_2} P_2''$, avec $u : P_1'' \to P_2''$}

\ill{} ouvert et fermé, contient $\xi'$, est stable par $\pi$, donc lui
correspond un \uncertain{sous-préschéma} ouvert et fermé $S'$ de $P'/\pi$ :
le \ill{} ouvert et fermé \marginal{ensemblistement} [donc
\uncertain{fidèlement} plat et quasi-compact \ill{}]. De plus, on voit
\uncertain{que} les \struck{\ill{}} fibrés géom. de $S' \to S$ \ill{}
portent un tel ; donc \ill{}

\page{166}

Supposons \struck{$S$} $S$ loc. noeth. \emph{ou} $\pi$ artinien (condition des
chaînes descendantes)].

\begin{itemize}
\item[d)] \ill{} $c \in \varphi(\pi)$ \ill{} provient \ill{} d'un élément
\uncertain{minimal} $c' \in \varphi(\pi')$.
\end{itemize}

En effet, si $c$ correspond à $(P,\xi)$, \uncertain{considérons} \ill{} et à
$(P'', y'')$ [pour $\pi'' \subset \pi$]. Considérons
\[
  P/\pi' \cap \pi'' \longrightarrow P/\pi' \times P/\pi''
\]
\marginal{$\pi/\pi' \cap \pi'' \to \pi/\pi' \times \pi/\pi''$}
c'est une immersion ouverte et fermée, il \ill{} la deuxième \ill{}
\ill{} un pt $(y',y'')$ est dans l'image, \ill{}

\ill{} en particulier \ill{}. Donc on peut restreindre le groupe structural
à $\pi' \cap \pi''$. Donc les groupes auxquels on peut restreindre forment un
\ill{} filtrant \emph{décroissant} $\pi_i$, \ill{} des $P_i \subset P$
\ill{}

\struck{\ill{}}

\note{le dernier tiers de la page est biffé de plusieurs traits diagonaux ;
rien ne s'en laisse lire}

\page{167}

principal homogène sous $\pi' = \bigcap \pi_i$. On voit ici que la question est
locale sur $S$ [et \ill{} \uncertain{fini}] donc on \uncertain{peut supposer}
$S$ \ill{}. Mais \ill{}

\marginal{$P^0 \leftarrow P$ au-dessus de $S' \to S$}

il faut $S^{(0)}$ de type fini sur $\mathbb{Z}$, et $(P^{(0)}, \xi^{(0)})$
principal homogène sous \ill{} qui donne $P$ par extension \ill{} [et
\ill{} par $P$ est trivialisé par $S' \to S$ \ill{}]. On voit qu'on
\uncertain{peut supposer} $S$ \uncertain{noethérien}.

\struck{Alors les} l'on tire \uncertain{un meilleur résultat} : on regarde le
\uncertain{composant} connexe $P^0$ de $\xi$ dans $P$ [\struck{exact}
\uncertain{qui est ouvert} et \emph{fermé}, car $P$ est loc.
\uncertain{noeth.} donc loc. connexe]. Je dis que $P^0 \to S$ est
\emph{surjectif}. Comme l'image \ill{}, \uncertain{plus généralement} elle
est fermée, on est ramené au cas d'un \uncertain{anneau de valuation},
\uncertain{disons} \uncertain{complet} : \ill{} que si $P$ est \ill{}
\ill{} sur $S$

\page{168}

\emph{Application.} Une relation étale est dite à \emph{groupe structural}
$\pi$, \struck{$\pi$, faisceau $F$, \ill{}} \ill{} \struck{dans un schéma en
$\underline{\mathrm{Aut}}_S(T)$}, qui \ill{} donne un \uncertain{revêtement
principal} de $T$ \ill{} et une \ill{} $\pi_S$ \uncertain{principal} de
\ill{}
\[
  \underline{\mathrm{Isom}}_S(E_S, T).
\]
Si $S$ est \uncertain{connexe} \uncertain{muni d'un} pt géom. $\xi$, il vient
\ill{} $=$ \ill{} \uncertain{donner} un $T$ \ill{} \ill{} telle
structure, ou de se donner des opérations \struck{continues} de
$\widetilde{\pi}_1(S,\xi)$ \struck{\ill{}} sur $F$ !

\ill{} donner $\widetilde{\pi}_1(S,\rho) \longrightarrow \struck{\ill{}}\pi$.

\emph{Exemple} \uncertain{Schémas en groupes commutatifs}, \ill{} torsion sur
$S$, \struck{à fibres \ill{}} \uncertain{à génération} finie sur $S$,
correspondant aux $(\widetilde{\pi}_1(S,\xi))$-Modules \uncertain{admissibles}
\struck{de présentation} qui sont des $\mathbb{Z}$-modules de type fini
\ill{}

\page{169}

\ill{} de groupes \ill{} \emph{$\pi_S$}, donc \ill{} la \emph{composante
connexe} $P^0$ est surjective sur $S$. \ill{}, \ill{} $P$ est
\emph{trivial} !

En effet, comme $P^{(0)} \subset \bigcap P^{(i)}$ \ill{}, on voit que
$\bigcap$ pris \ill{} \ill{} de gpes fin\ill{} \struck{\ill{}}
\struck{finie} \marginal{chacun} des $P^{(i)}$ est une \uncertain{réunion} de
\uncertain{composantes connexes} de $P$, il \ill{} de l'\uncertain{intersection},
\ill{} $P'$, \ill{} $P^0$.

Elle est \emph{stable} sous $\pi'$, et \ill{} \uncertain{ensemble} (car
\ill{} \ill{} fibres \uncertain{géométriques}, et \ill{} que
$P_s' \neq \emptyset$]. \ill{} $P'$ est \ill{} fibré principal homogène
sous $\pi'$. cqfd.

\marginal{\ill{} $\widetilde{\pi}_1(S,\xi)$ \ill{}}

\emph{Conséquence} Si $\pi$ varie dans \struck{la} \uncertain{une}
\uncertain{catégorie} \ill{} des gpes artiniens, \struck{\ill{}}
\uncertain{stable} par \uncertain{produits finis} et \uncertain{sous-groupes},
\ill{} si $\pi$ varie \ill{} \ill{} \uncertain{opérations}, \ill{}
\ill{} le foncteur
\[
  \pi \longmapsto \pi^1(S,\xi;\pi)
\]
est \emph{strictement pro-représentable}.

\section*{Groupe fondamental : revêtements infinis ; les cinq classes $C_1, \ldots, C_5$ (pages 170 et 171)}

\page{170}

\emph{Groupe fondamental} : (Revêtements infinis)

$S$ préschéma. Pour tout groupe discret $G$, on définit
\[
  H^1_{\mathcal{C}}(S,G) = \varinjlim_{T/S} H^1(T/S, G)
\]
\note{la ligne qui suit « $T/S$ » sous le signe de limite est biffée et ne se
laisse pas lire}
où
\[
  H^1(T/S, G) = H^1(K_{T/S}, G) = H^1(\pi_0(K_{T/S}), G)
\]
et où $\mathcal{C}$ est une catégorie de $S$-préschémas.

\marginal{élucider ici les questions de \uncertain{fonctorialité} \ill{}}

Les cas les plus importants sont les suivants :

\begin{itemize}
\item[(i)] $C_1$ \quad $S$-préschémas fid. plats quasi-compacts
\struck{\ill{}}.
\item[(ii)] $C_2$ \quad ————— de type fini
\item[(iii)] $C_3$ \quad ————— quasi-finis
\item[(iv)] $C_4$ \quad ————— finis \uncertain{principaux}
\item[(v)] $C_5$ \quad ————— \struck{étales} finis, étales
\uncertain{principaux}
\end{itemize}

On a donc $H^1_{C_i}(\struck{P}S, G) \hookrightarrow H^1_{C_j}(S,G)$ \ill{}
application injective canonique. Si $G$ est \emph{fini}, alors \ill{}, les
hom. sont bijectifs, et $H^1(S,G)$ est l'ensemble des \uncertain{classes}
(à isom. près) \uncertain{de} revêtements étales principaux de groupe $G$,
donc $\simeq \mathrm{Hom}(\pi_1(S,a), G)/\mathrm{int}(G)$. Notre
\uncertain{intérêt} ici \uncertain{est} \ill{} le cas $G$ \emph{infini}.
Notons en passant
\[
  H^1_{C_3} \simeq H^1_{C_2}
\]

\page{171}

et si $S$ est le spectre d'un anneau de Dedekind
\[
  H^1_{C_2} \xrightarrow{\ \sim\ } H^1_{C_1}
\]
[d'où dans ce cas $H^1_{C_3} \simeq H^1_{C_2} \simeq H^1_{C_1}$]. Si $S$ est le
spectre d'un anneau \struck{local complet} \marginal{de valuation} discrète à
corps résiduel alg. clos, on aura de \uncertain{même}
\[
  H^1_{C_i} \simeq H^1(\widehat{V}/V, -) \qquad \text{pour } i = 1,2,3
\]
\note{entre $H^1_{C_i}$ et le signe $\simeq$ la page porte un caractère isolé
qui ne se laisse pas identifier ; l'indice du $H^1$ de droite est biffé}

Si $S$ est le spectre d'un anneau local \emph{complet} $V$, alors
\struck{\ill{}}
\[
  H^1_{C_5} \simeq H^1_{C_4} \simeq H^1_{C_3} \quad (\simeq H^1_{C_2})
\]
et on aura de plus
\[
  H^1_{C_i}(V,-) \simeq H^1_{C_i}(k, -) \qquad (i = 2,3,4,5)
\]
\note{sous le $k$ la page porte « corps résiduel »}

\section*{Plan en sept sections pour les deux théorèmes de changement de base (page 173)}

\note{le feuillet qui ouvre la suite --- page 172, par ailleurs nu --- porte de
sa main, en haut à droite : « Plan »}

\page{173}

\emph{Groupe fondamental}

\begin{itemize}
\item[I] \struck{Théorème de comparaison \ill{}} \marginal{Premier} th. de
changement de base pour $\pi_1$ (ou th. de comparaison)
\item[II] Deuxième th. de changement de base pour $\pi_1$
\end{itemize}

\begin{enumerate}
\item Morphismes $(-1)$-connexes et $0$-connexes
\item Morphismes locaux $(-1)$-connexes et $0$-connexes.
\item \struck{Une} Caractérisation des morphismes locaux ou globaux
$i$-connexes ($i = 0,1$).
\item \struck{Le} Deuxième théorème de \struck{comparaison}
\marginal{du chgt de base pour $\pi_1$ :} énoncés et corollaires \ldots
\item \struck{\ill{}}, \struck{et} Démonstration des théorèmes.
\item Généralisation à l'aide de la résolution des singularités.
\item Application : la dimension cohomologique de certains corps \ldots
\end{enumerate}

\section*{$\pi_2$ : la suite exacte de Hochschild--Serre contre la tour des revêtements (pages 175 et 176)}

\note{le feuillet qui ouvre la suite --- page 174, par ailleurs nu --- porte de
sa main, en haut à droite : « $\Pi_2$ »}

\page{175}

\begin{tikzcd}[column sep=small, row sep=small, nodes={font=\scriptsize}]
\widetilde{\widetilde{X}}{}^i \arrow[d] & \\
\widetilde{X}^i \arrow[d] & \widetilde{X} \arrow[d, no head] \\
X^i \arrow[r] & X
\end{tikzcd}

\note{sur la page la colonne de gauche est prise dans une accolade portant
$\widetilde{\pi}$ à gauche et, à droite, une mention à deux étages où se lisent
$\widehat{\widetilde{\pi}}$ et $\widetilde{\pi} = \pi$}

\[
  \pi^{\mathrm{I}}(X,G) \simeq H^2(\check{X},G) \simeq \varinjlim H^2(X_i,G)
  \simeq \varinjlim H^2(\widetilde{X}_i,G)
\]
\note{l'exposant du premier terme est tracé comme un « I » ; plus bas sur la
même page il écrit $\pi^2$}

\marginal{$\widetilde{\pi}/H_i = \pi_i$}

\[
  H^*(\widetilde{X}_i,G) \Longleftarrow E_2^{p,q} = H^p(H_i, H^q(\widetilde{X},G))
\]
\[
  \struck{\ill{}} H^*(\widetilde{X},G) \Longleftarrow E^{p,q} =
  \varinjlim_i H^p(H_i, H^q(\widetilde{X},G))
\]

\[
  \cdots \to \varinjlim H^2(H_i,G) \to \pi^2(X,G) \to
  \varinjlim_i \pi^2(\widetilde{X}_i,G)^{H_i} \to \varinjlim H^3(H_i,G)
\]
\note{sous le troisième terme la page porte $\pi^2(X^k,G)$, sous le quatrième
$H^3(\widetilde{X},G)$ ; à gauche et à droite deux termes sont entourés, et
au-dessus de la suite est esquissé le premier quadrant $(p,q)$ de la suite
spectrale}

\marginal{$\pi_2(X) \to G$}

\[
  H^*(X^k,G) \Longleftarrow H
\]
\[
  \mathrm{Ker}\,\bigl(\pi^2(X,G) \to \pi^2(X^k,G)\bigr) \simeq
  \varinjlim_i H^2(H_i,G)
\]
\[
  \mathrm{Coker}\,(\text{id}) \simeq \mathrm{Ker}\Bigl(\varinjlim H^3(H_i,G)
  \longrightarrow H^3(\widetilde{X},G)\Bigr)
\]

\page{176}

\begin{tikzcd}[column sep=small, row sep=small, nodes={font=\scriptsize}]
X \arrow[d] & X' \arrow[l] \arrow[d] & X'' \arrow[l, bend left=12]
  \arrow[l, bend right=12] \arrow[d] \\
Y & Y' \arrow[l] & Y'' \arrow[l, bend left=12] \arrow[l, bend right=12]
\end{tikzcd}

\note{au-dessus de $X'$ et de $X''$ la page porte $Z'$ et
$Z_1'' \simeq Z_2''$ ; sous $Y'$ un double trait la relie à $X$, et sous $Y''$
une flèche monte de $\widetilde{X}''$}

\[
  \pi_1(\overline{F}) \to \pi_1(X) \to \pi_1(Y) \to e
\]
\note{de $\pi_1(\overline{F})$ part une flèche vers $G$, et de $\pi_1(X)$ une
flèche pointillée vers le même $G$}

\[
  \pi_1(Y) \longrightarrow G^{\pi_1(X)}
\]
\[
  H^1(\pi_1(X), G) \longrightarrow H^1(\pi_1(\overline{F}), G)^{\pi_1(Y)}
\]
\[
  H^1\bigl(\pi_1(Y), G^{\pi_1(\overline{F})}\bigr) =
  H^1(\pi_1(Y), \mathfrak{z})
\]
\note{la flèche va du second $H^1$ vers le premier}

\[
  \pi_2(Y) \longrightarrow \mathfrak{z} \qquad H^1(\pi_1(\overline{F}), G)
\]
\[
  \pi_1(Y') \longrightarrow \mathfrak{z}
\]
\[
  H^1\bigl(Y'/Y, \mathrm{Hom}(\pi_1(\overline{F}), \mathfrak{z})\bigr) = 0
\]
\note{au-dessus de $\mathrm{Hom}$ la page porte un $H^1$ biffé}

\begin{tikzcd}[column sep=small, row sep=small, nodes={font=\scriptsize}]
\pi_1(Y) & \pi_1(Y') \arrow[l] & \pi_1(Y'') \arrow[l, bend left=12]
  \arrow[l, bend right=12] \arrow[d] & \pi_1(Y''') \arrow[l, bend left=18]
  \arrow[l] \arrow[l, bend right=18] \\
& & \mathfrak{z} &
\end{tikzcd}

\section*{$\pi_0$, $\pi_1$ et formalisme simplicial (revêtements infinis) (pages 178 à 180)}

\note{le feuillet qui ouvre la suite --- page 177, par ailleurs nu --- porte de
sa main, en haut à droite : « $\pi_0$, $\pi_1$ et formalisme simplicial
(Revêtements infinis) ». Le titre de cette section en vient}

\page{178}

\emph{Proposition} $S$ un préschéma \emph{connexe}, \struck{\ill{}}
$S'$ un préschéma sur $S$ tel que $S' \to S$ soit \emph{submersif}
[p. ex. fid. plat et quasi-compact, ou propre surjectif], $S''$, $S'''$ les
carrés, cubes fibrés.

On suppose que $S'$ est \emph{somme} de préschémas \emph{connexes} (i.e. ses
composantes connexes sont ouvertes, i.e. $\pi_0(S')$ discret)
\struck{et l'ens. des composantes connexes de $S'$ \ill{}}

Soit $\xi$ un pt [géométrique] de $S$. Alors pour un groupe \struck{discret}
$G$ variable, le foncteur
\[
  G \rightsquigarrow \pi^1(S'/S, \xi; G)
\]
\struck{contravariant} [où $\pi^1(S'/S,\xi;G)$ désigne \marginal{l'ens. des
classes} \ill{} \uncertain{à isom.} près, de données de descente
\uncertain{sur $G_{S'}$} relatives à $S' \to S$, ponctuées au $\xi$] est
strict. pro-représentable :
\[
  \pi^1(S'/S, \xi; G) = \varinjlim \mathrm{Hom}(\pi_i, G)
\]

Si $\pi_0(S'')$ est \marginal{quasi-}\emph{compact}, i.e. les fonctions loc.
constantes sur $\pi_0(S'')$ \ill{} \uncertain{à valeurs finies} dans des
ensembles, --- p. ex. si $S''$ quasi-compact] \struck{\ill{}}
\struck{\ill{}} alors les $\pi_i$ sont de type fini.

Si $S''$ a un nb fini \marginal{$''$} de composantes connexes, alors le
foncteur précédent est \uncertain{représentable} : \struck{\ill{}}
\[
  \pi^1(S'/S, \xi; G) \simeq \mathrm{Hom}(\pi_1(S'/S, \xi), G)
\]
avec $\pi_1(S'/S,\xi)$ \marginal{groupe discret} \struck{\ill{}} de type
fini. \ill{} si \ill{} des composantes connexes de $S'$ est
\uncertain{finie}, et \ill{} fini, et \struck{\ill{}} alors
$\pi_1(S'/S;\xi)$ est \uncertain{engendré} par $n''$ \uncertain{générateurs}.

\page{179}

\emph{Enfin}, considérons le \struck{\ill{}} foncteur
\[
  G \rightsquigarrow \pi^1_{\mathrm{eff}}(S'/S, \xi; G)
\]
correspondant aux \struck{\ill{}} données de descente \emph{effectives}
ponctuées sur $\xi$. \struck{\ill{}}

Ce foncteur est également \emph{pro-représentable} [par un système projectif de
quotients des $\pi_i$].

\emph{N.B.} J'ignore, --- si $S,S',S''$ noethériens, \ill{} dernier foncteur
est $=$ \emph{représentable}, à fortiori (comme $\pi^1(S'/S,\xi;G)$ est repr.
dans ce cas) j'ignore si dans ce cas $\pi^1 = \pi^1_{\mathrm{eff}}$, i.e.
\emph{toute} donnée de descente sur $G_{S'}$ est \emph{effective} \ill{}
C'est vrai cependant \struck{\ill{}} \uncertain{hypothèse} sur $S''$, $S$ et
$S'$ \ill{} et si loc. noeth.] si \struck{\ill{}} dans $G_i$ \ill{}
intersection de sous-gpes \ill{} d'indice fini \ldots

\marginal{Pb \uncertain{intér.} pour l'affirmative [ \ill{} et
$S' \to S$ \uncertain{fid. plat quasi-compact} ]}

\page{180}

\begin{tikzcd}[column sep=small, row sep=small, nodes={font=\scriptsize}]
Y & X_0^0 \arrow[l, dashed] & X_1^1 \arrow[l, bend left=12]
  \arrow[l, bend right=12] & X_2^2 \arrow[l, bend left=20] \arrow[l]
  \arrow[l, bend right=20] \\
\xi \arrow[u] & X_0^{\xi} \arrow[l] \arrow[u] & X_1^{\xi}
  \arrow[l, bend left=12] \arrow[l, bend right=12] \arrow[u] & X_2^{\xi}
  \arrow[l, bend left=20] \arrow[l] \arrow[l, bend right=20] \arrow[u] \\
& K_0^0 & K_1^1 \arrow[l, bend left=12] \arrow[l, bend right=12] & K_2^2
  \arrow[l, bend left=20] \arrow[l] \arrow[l, bend right=20] \\
& k_0^0 \arrow[u, "i_0"] & k_1^1 \arrow[l] \arrow[u, "i_1"] & k_2^2
  \arrow[l, bend left=12] \arrow[l, bend right=12] \arrow[u, "i_2"]
\end{tikzcd}

\[
  K_i^{\xi} = \pi_0(X_i^{\xi})
\]

$\pi_0(K) = {}$ \ill{} \ldots $K^0 = e$, et au-dessous
$\pi_0(k) = {}$ « \ill{} » \ldots $k^0 = e$.

\note{les deux lignes sont reliées par une flèche montante ; ce qui sépare
chaque signe d'égalité de sa conclusion ne se laisse pas lire}

Données de descente sur $E_{X_0}$ $\simeq$ Données de descente sur $E_{K_0}$
\uncertain{relatives} à $(\cdots)$

$\simeq$ applications \struck{\ill{}} ensemblistes $f : K_1 \to G$ telles que
\[
  f(p_1(\xi)) = f(p_0(q_2))\, f(p_2(q_1)) \qquad \text{pour tt } \xi \in K
\]
Équivalence de telles :
\[
  h(p_1(\xi))\, f(\xi)\, h(p_0 q_1)^{-1} \simeq f(\xi)
\]

\emph{Données de descente ponctuées} sur $E_{K_0}$

$\simeq$ applications ensemblistes $f : K_1 \to G$ telles que
\[
  \begin{cases}
    f(i_1(q)) = \struck{1}\, e_G & \text{pour tt } q \in k_1 \\
    \struck{\ill{}}\quad f \circ p_1 = (f \circ p_0)(f \circ p_2) &
  \end{cases}
\]
Équivalence de telles :
\[
  h(p_1 q)\, f(q)\, h(p_0 q)^{-1} \simeq f(q)
\]
\uncertain{Si} $h(i_0(q)) = 1$ pour tt $i_0$.

\end{document}
